\documentclass[11pt,a4paper]{article}

\usepackage[utf8]{inputenc}
\usepackage[T1]{fontenc}
\usepackage[french]{babel}

\usepackage{styles/couleurs}
\usepackage{styles/mise_en_page}
\usepackage{styles/boites}
\usepackage{styles/maths}
\usepackage{styles/figures}

\usepackage{amsmath}
\usepackage{amssymb}
\usepackage[most]{tcolorbox}
\usepackage{tikz}
\usetikzlibrary{calc,patterns}

% =========================================================
% MISE EN PAGE
% =========================================================

\geometry{
  top=5mm,
  bottom=5mm,
  left=9mm,
  right=9mm
}

\pagestyle{empty}

\setlength{\parindent}{0pt}
\setlength{\parskip}{0pt}

\setlength{\abovedisplayskip}{3pt}
\setlength{\belowdisplayskip}{3pt}
\setlength{\abovedisplayshortskip}{2pt}
\setlength{\belowdisplayshortskip}{2pt}

% =========================================================
% TITRE FICHE ÉLÈVE
% =========================================================

\newcommand{\titreprobleme}[2]{%
\noindent
\begin{minipage}{\linewidth}
\centering
{\sffamily\LARGE\bfseries\color{TitreC}
Résoudre un problème avec Pythagore}

\vspace{0.5mm}

{\sffamily\large\bfseries\color{SavoirC}
Problème #1 -- #2}
\end{minipage}

\vspace{1mm}
}

% =========================================================
% TITRE CORRECTION
% =========================================================

\newcommand{\titrecorrection}[1]{%
\noindent
\begin{minipage}{\linewidth}
\centering
{\sffamily\LARGE\bfseries\color{TitreC}
Un exemple de correction complète}

\vspace{0.5mm}

{\sffamily\large\bfseries\color{SavoirC}
#1}
\end{minipage}

\vspace{2mm}
}

% =========================================================
% BOÎTES ÉLÈVE
% =========================================================

\newtcolorbox{enoncebox}{
  enhanced,
  colback=white,
  colframe=ActionC!55,
  boxrule=0.7pt,
  arc=1.5mm,
  left=4mm,
  right=4mm,
  top=1.5mm,
  bottom=1.5mm,
  height=4cm,
  valign=center,
  before skip=0pt,
  after skip=0pt
}

\newtcolorbox{connaissancesbox}{
  enhanced,
  colback=SavoirBg,
  colframe=SavoirC!55,
  boxrule=0.7pt,
  arc=1.5mm,
  left=4mm,
  right=4mm,
  top=1.5mm,
  bottom=1.2mm,
  height=4.2cm,
  before skip=0pt,
  after skip=0pt
}

% =========================================================
% CASE À COCHER
% =========================================================

\newcommand{\casecocher}[1]{%
\noindent
\(\square\)\hspace{0.35em}%
\parbox[t]{0.89\linewidth}{#1}
\par
\vspace{0.15em}
}

% =========================================================
% CADRES À POINTS
%
% #1 largeur
% #2 hauteur
% #3 titre
% #4 couleur du titre
% =========================================================

\newcommand{\zonepoints}[4]{%
\begin{tikzpicture}[x=1cm,y=1cm,baseline=(current bounding box.north)]

  \fill[white]
    (0,0) rectangle (#1,#2);

  \fill[
    pattern=dots,
    pattern color=gray!25
  ]
    (0,0) rectangle (#1,#2);

  \draw[
    black!35,
    line width=0.6pt,
    rounded corners=1.5mm
  ]
    (0,0) rectangle (#1,#2);

  \node[
    anchor=north west,
    font=\sffamily\bfseries,
    text=#4,
    fill=white,
    inner xsep=3pt,
    inner ysep=2pt
  ]
  at (0.25,#2-0.18)
  {#3};

\end{tikzpicture}%
}

% =========================================================
% TROIS ZONES DE TRAVAIL
% =========================================================

\newcommand{\zonestravail}{%

\noindent
\begin{minipage}[t]{0.49\linewidth}
\vspace{0pt}
\centering

\zonepoints{9.35}{9}{Brouillon}{black!60}

\end{minipage}
\hfill
\begin{minipage}[t]{0.49\linewidth}
\vspace{0pt}
\centering

\zonepoints{9.35}{9}{Ma version}{ActionC}

\end{minipage}

\vspace{0.8mm}

\noindent
\makebox[\linewidth][c]{%
\zonepoints{19.15}{8.6}{Notre version}{SavoirC}%
}%
}

% =========================================================
% PAGE DE CORRECTION
% =========================================================

\newtcolorbox{correctionbox}{
  enhanced,
  colback=SavoirBg!30,
  colframe=SavoirC!40,
  boxrule=0.6pt,
  arc=2mm,
  left=7mm,
  right=7mm,
  top=5mm,
  bottom=5mm,
  height=26.2cm,
  valign=top,
  before skip=0pt,
  after skip=0pt
}

\newcommand{\pagecorrection}[2]{%

\titrecorrection{#1}

\begin{correctionbox}

#2

\end{correctionbox}
}

% =========================================================
% DOCUMENT
% =========================================================

\begin{document}

% =========================================================
% PROBLÈME 1 : LOSANGE
% =========================================================

\titreprobleme{1}{Le losange}

\begin{enoncebox}

\begin{minipage}[c]{0.64\linewidth}

Un jardin a la forme d'un losange \(ABCD\).

Chaque côté mesure \(27{,}4\) m et la diagonale \([AC]\)
mesure \(42\) m.

\smallskip

\textbf{Quelle est la longueur de la diagonale \([BD]\) ?}

\end{minipage}
\hfill
\begin{minipage}[c]{0.31\linewidth}

\centering

\begin{tikzpicture}[scale=0.57]

\coordinate (A) at (-2.5,0);
\coordinate (C) at (2.5,0);
\coordinate (B) at (0,2.1);
\coordinate (D) at (0,-2.1);

\draw[thick] (A)--(B)--(C)--(D)--cycle;
\draw[dashed] (A)--(C);
\draw[dashed] (B)--(D);

\node[left] at (A) {$A$};
\node[above] at (B) {$B$};
\node[right] at (C) {$C$};
\node[below] at (D) {$D$};

\node[above left] at ($(A)!0.5!(B)$)
{$27{,}4$ m};

\node[below] at ($(A)!0.5!(C)$)
{$42$ m};

\end{tikzpicture}

\end{minipage}

\end{enoncebox}

\vspace{0.8mm}

\zonestravail

\vspace{0.8mm}

\begin{connaissancesbox}

{\sffamily\bfseries\color{SavoirC}
Point connaissances -- à consulter si besoin}

\smallskip

Parmi ces connaissances, lesquelles peuvent t'aider ?
Coche celles que tu utilises.

\smallskip

\begin{minipage}[t]{0.49\linewidth}

\casecocher{Les diagonales d'un losange se coupent en leur milieu.}

\casecocher{Les diagonales d'un losange sont perpendiculaires.}

\casecocher{Un losange a quatre côtés de même longueur.}

\end{minipage}
\hfill
\begin{minipage}[t]{0.49\linewidth}

\casecocher{Le périmètre d'un losange vaut quatre fois la longueur d'un côté.}

\casecocher{L'aire d'un triangle vaut
\(\dfrac{\text{base}\times\text{hauteur}}{2}\).}

\casecocher{Dans un triangle rectangle, on peut utiliser
le théorème de Pythagore.}

\end{minipage}

\end{connaissancesbox}

\newpage

% =========================================================
% CORRECTION 1
% =========================================================

\pagecorrection{Problème 1 -- Le losange}{

Les diagonales d'un losange sont perpendiculaires et se coupent
en leur milieu.

On note \(O\) leur point d'intersection.

Comme \(O\) est le milieu de \([AC]\) :

\[
AO=\frac{AC}{2}
=\frac{42}{2}
=21\text{ m}.
\]

Le triangle \(AOB\) est rectangle en \(O\).

D'après le théorème de Pythagore :

\[
AB^2=AO^2+OB^2
\]

\[
27{,}4^2=21^2+OB^2
\]

\[
750{,}76=441+OB^2
\]

\[
OB^2=750{,}76-441
\]

\[
OB^2=309{,}76
\]

\[
OB=\sqrt{309{,}76}
\]

\[
OB=17{,}6\text{ m}.
\]

Comme \(O\) est aussi le milieu de \([BD]\) :

\[
BD=2\times OB
\]

\[
BD=2\times17{,}6
\]

\[
\boxed{BD=35{,}2\text{ m}}
\]

La seconde diagonale du losange mesure donc
\(\boxed{35{,}2\text{ m}}\).
}

\newpage

% =========================================================
% PROBLÈME 2 : TRAPÈZE
% =========================================================

\titreprobleme{2}{Le terrain à clôturer}

\begin{enoncebox}

\begin{minipage}[c]{0.62\linewidth}

Un terrain a la forme du trapèze rectangle \(ABCD\).

\[
AB=9\text{ m},
\qquad
AD=12\text{ m},
\qquad
DC=14\text{ m}.
\]

On souhaite poser une clôture tout autour du terrain.

\smallskip

\textbf{Quelle longueur de grillage faut-il prévoir ?}

\end{minipage}
\hfill
\begin{minipage}[c]{0.32\linewidth}

\centering

\begin{tikzpicture}[scale=0.54]

\coordinate (D) at (0,0);
\coordinate (A) at (0,3);
\coordinate (B) at (3.3,3);
\coordinate (C) at (5.1,0);

\draw[thick] (A)--(B)--(C)--(D)--cycle;

\draw[thick]
(D) ++(0,0.35)
-- ++(0.35,0)
-- ++(0,-0.35);

\draw[thick]
(A) ++(0,-0.35)
-- ++(0.35,0)
-- ++(0,0.35);

\node[above left] at (A) {$A$};
\node[above right] at (B) {$B$};
\node[below right] at (C) {$C$};
\node[below left] at (D) {$D$};

\node[above] at ($(A)!0.5!(B)$)
{$9$ m};

\node[left] at ($(A)!0.5!(D)$)
{$12$ m};

\node[below] at ($(D)!0.5!(C)$)
{$14$ m};

\end{tikzpicture}

\end{minipage}

\end{enoncebox}

\vspace{0.8mm}

\zonestravail

\vspace{0.8mm}

\begin{connaissancesbox}

{\sffamily\bfseries\color{SavoirC}
Point connaissances -- à consulter si besoin}

\smallskip

Parmi ces connaissances, lesquelles peuvent t'aider ?
Coche celles que tu utilises.

\smallskip

\begin{minipage}[t]{0.49\linewidth}

\casecocher{Le périmètre d'une figure est la somme
des longueurs de ses côtés.}

\casecocher{Les côtés opposés d'un rectangle ont la même longueur.}

\casecocher{Dans un triangle rectangle, on peut utiliser
le théorème de Pythagore.}

\end{minipage}
\hfill
\begin{minipage}[t]{0.49\linewidth}

\casecocher{L'aire d'un trapèze dépend de ses deux bases
et de sa hauteur.}

\casecocher{Une longueur peut parfois être obtenue par différence.}

\casecocher{Un rectangle possède quatre angles droits.}

\end{minipage}

\end{connaissancesbox}

\newpage

% =========================================================
% CORRECTION 2
% =========================================================

\pagecorrection{Problème 2 -- Le terrain à clôturer}{

On trace la perpendiculaire à \([DC]\) passant par \(B\).
Elle coupe \([DC]\) en \(H\).

Le quadrilatère \(ABHD\) est un rectangle.

On a donc :

\[
DH=AB=9\text{ m}
\]

et

\[
BH=AD=12\text{ m}.
\]

Ainsi :

\[
HC=DC-DH
\]

\[
HC=14-9
\]

\[
HC=5\text{ m}.
\]

Le triangle \(BHC\) est rectangle en \(H\).

D'après le théorème de Pythagore :

\[
BC^2=BH^2+HC^2
\]

\[
BC^2=12^2+5^2
\]

\[
BC^2=144+25
\]

\[
BC^2=169
\]

\[
BC=\sqrt{169}
\]

\[
BC=13\text{ m}.
\]

Le périmètre du terrain vaut :

\[
AB+BC+CD+DA
\]

\[
9+13+14+12=48\text{ m}.
\]

\[
\boxed{\text{Il faut prévoir }48\text{ m de grillage}.}
\]
}

\newpage

% =========================================================
% PROBLÈME 3 : TREMPLIN
% =========================================================

\titreprobleme{3}{Le tremplin}

\begin{enoncebox}

\begin{minipage}[c]{0.63\linewidth}

La surface inclinée d'un tremplin est un rectangle
de largeur \(0{,}8\) m.

Le tremplin s'élève de \(0{,}9\) m sur une longueur horizontale
de \(1{,}2\) m.

Le revêtement de la surface inclinée a coûté \(128{,}52\) €.

\smallskip

\textbf{Quel est le prix du revêtement au mètre carré ?}

\end{minipage}
\hfill
\begin{minipage}[c]{0.31\linewidth}

\centering

\begin{tikzpicture}[scale=0.74]

\coordinate (A) at (0,0);
\coordinate (B) at (3.4,0);
\coordinate (C) at (3.4,2.1);

\draw[thick] (A)--(B)--(C)--cycle;

\draw[thick]
(B) ++(-0.28,0)
-- ++(0,0.28)
-- ++(0.28,0);

\node[below left] at (A) {$A$};
\node[below right] at (B) {$B$};
\node[above right] at (C) {$C$};

\node[below] at ($(A)!0.5!(B)$)
{$1{,}2$ m};

\node[right] at ($(B)!0.5!(C)$)
{$0{,}9$ m};

\node[above left] at ($(A)!0.5!(C)$)
{$?$};

\end{tikzpicture}

\smallskip

Largeur : \(0{,}8\) m

\end{minipage}

\end{enoncebox}

\vspace{0.8mm}

\zonestravail

\vspace{0.8mm}

\begin{connaissancesbox}

{\sffamily\bfseries\color{SavoirC}
Point connaissances -- à consulter si besoin}

\smallskip

Parmi ces connaissances, lesquelles peuvent t'aider ?
Coche celles que tu utilises.

\smallskip

\begin{minipage}[t]{0.49\linewidth}

\casecocher{L'aire d'un rectangle vaut
longueur \(\times\) largeur.}

\casecocher{Le volume d'un prisme droit vaut
aire de la base \(\times\) hauteur.}

\casecocher{Dans un triangle rectangle,
on peut utiliser le théorème de Pythagore.}

\end{minipage}
\hfill
\begin{minipage}[t]{0.49\linewidth}

\casecocher{Un prix au mètre carré peut être obtenu
en divisant le prix total par l'aire.}

\casecocher{Le périmètre d'un rectangle dépend
de sa longueur et de sa largeur.}

\casecocher{Deux droites perpendiculaires forment un angle droit.}

\end{minipage}

\end{connaissancesbox}

\newpage

% =========================================================
% CORRECTION 3
% =========================================================

\pagecorrection{Problème 3 -- Le tremplin}{

Le triangle \(ABC\) est rectangle en \(B\).

D'après le théorème de Pythagore :

\[
AC^2=AB^2+BC^2
\]

\[
AC^2=1{,}2^2+0{,}9^2
\]

\[
AC^2=1{,}44+0{,}81
\]

\[
AC^2=2{,}25
\]

\[
AC=\sqrt{2{,}25}
\]

\[
AC=1{,}5\text{ m}.
\]

La surface inclinée du tremplin est un rectangle
de longueur \(1{,}5\) m et de largeur \(0{,}8\) m.

Son aire vaut :

\[
1{,}5\times0{,}8
=1{,}2\text{ m}^2.
\]

Le prix du revêtement pour \(1\text{ m}^2\) vaut :

\[
128{,}52\div1{,}2
=107{,}1.
\]

\[
\boxed{\text{Le revêtement coûte }107{,}10\text{ € par m}^2.}
\]
}

\newpage

% =========================================================
% PROBLÈME 4 : PARC
% =========================================================

\titreprobleme{4}{Le raccourci dans le parc}

\begin{enoncebox}

\begin{minipage}[c]{0.64\linewidth}

Un parc rectangulaire mesure \(60\) m de large et \(80\) m de long.

Pour aller de \(A\) à \(C\), on peut suivre les côtés \(AB\) puis \(BC\),
ou traverser directement le parc par la diagonale \([AC]\).

\smallskip

\textbf{Combien de mètres gagne-t-on en prenant le raccourci ?}

\end{minipage}
\hfill
\begin{minipage}[c]{0.30\linewidth}

\centering

\begin{tikzpicture}[scale=0.53]

\coordinate (A) at (0,0);
\coordinate (B) at (4,0);
\coordinate (C) at (4,3);
\coordinate (D) at (0,3);

\draw[thick] (A)--(B)--(C)--(D)--cycle;
\draw[dashed,thick] (A)--(C);

\node[below left] at (A) {$A$};
\node[below right] at (B) {$B$};
\node[above right] at (C) {$C$};
\node[above left] at (D) {$D$};

\node[below] at ($(A)!0.5!(B)$)
{$80$ m};

\node[right] at ($(B)!0.5!(C)$)
{$60$ m};

\end{tikzpicture}

\end{minipage}

\end{enoncebox}

\vspace{0.8mm}

\zonestravail

\vspace{0.8mm}

\begin{connaissancesbox}

{\sffamily\bfseries\color{SavoirC}
Point connaissances -- à consulter si besoin}

\smallskip

Parmi ces connaissances, lesquelles peuvent t'aider ?
Coche celles que tu utilises.

\smallskip

\begin{minipage}[t]{0.49\linewidth}

\casecocher{Un rectangle possède quatre angles droits.}

\casecocher{L'aire d'un rectangle vaut
longueur \(\times\) largeur.}

\casecocher{Dans un triangle rectangle,
on peut utiliser le théorème de Pythagore.}

\end{minipage}
\hfill
\begin{minipage}[t]{0.49\linewidth}

\casecocher{Pour connaître une distance gagnée,
on peut comparer deux trajets.}

\casecocher{Les diagonales d'un rectangle ont la même longueur.}

\casecocher{Le périmètre d'un rectangle dépend
de sa longueur et de sa largeur.}

\end{minipage}

\end{connaissancesbox}

\newpage

% =========================================================
% CORRECTION 4
% =========================================================

\pagecorrection{Problème 4 -- Le raccourci dans le parc}{

Le trajet qui suit les deux côtés mesure :

\[
AB+BC
\]

\[
80+60=140\text{ m}.
\]

Comme \(ABCD\) est un rectangle,
le triangle \(ABC\) est rectangle en \(B\).

D'après le théorème de Pythagore :

\[
AC^2=AB^2+BC^2
\]

\[
AC^2=80^2+60^2
\]

\[
AC^2=6400+3600
\]

\[
AC^2=10000
\]

\[
AC=\sqrt{10000}
\]

\[
AC=100\text{ m}.
\]

Le raccourci mesure donc \(100\) m.

La distance gagnée vaut :

\[
140-100=40\text{ m}.
\]

\[
\boxed{\text{On gagne }40\text{ m en prenant le raccourci}.}
\]
}

\newpage

% =========================================================
% PROBLÈME 5 : TOIT
% =========================================================

\titreprobleme{5}{Le toit de l'abri}

\begin{enoncebox}

\begin{minipage}[c]{0.62\linewidth}

La façade d'un abri mesure \(8\) m de large.

Le sommet \(S\) du toit est situé exactement au-dessus du milieu \(M\)
de la façade, à \(3\) m au-dessus du haut des murs.

Les plaques utilisées pour couvrir un versant mesurent \(4{,}8\) m.

\smallskip

\textbf{Une plaque est-elle assez longue pour couvrir un versant ?}

\end{minipage}
\hfill
\begin{minipage}[c]{0.33\linewidth}

\centering

\begin{tikzpicture}[scale=0.47]

% Façade
\coordinate (A) at (-3,0);
\coordinate (B) at (3,0);
\coordinate (C) at (3,-2.1);
\coordinate (D) at (-3,-2.1);

% Toit
\coordinate (M) at (0,0);
\coordinate (S) at (0,2.5);

% Murs
\draw[thick]
(A)--(B)--(C)--(D)--cycle;

% Toit
\draw[thick]
(A)--(S)--(B);

% Hauteur
\draw[dashed]
(S)--(M);

% Angle droit
\draw[thick]
($(M)+(-0.28,0)$)
-- ++(0,0.28)
-- ++(0.28,0);

% Milieu de AB
\draw[thick]
($(M)+(-0.13,-0.08)$)
--
($(M)+(0.13,0.08)$);

% Noms
\node[left] at (A) {$A$};
\node[right] at (B) {$B$};
\node[below right] at (M) {$M$};
\node[above] at (S) {$S$};

% Mesures
\node[below] at ($(A)!0.5!(B)+(0,-0.28)$)
{$8$ m};

\node[right] at ($(M)!0.5!(S)$)
{$3$ m};

\end{tikzpicture}

\end{minipage}

\end{enoncebox}

\vspace{0.8mm}

\zonestravail

\vspace{0.8mm}

\begin{connaissancesbox}

{\sffamily\bfseries\color{SavoirC}
Point connaissances -- à consulter si besoin}

\smallskip

Parmi ces connaissances, lesquelles peuvent t'aider ?
Coche celles que tu utilises.

\smallskip

\begin{minipage}[t]{0.49\linewidth}

\casecocher{Le milieu d'un segment partage ce segment
en deux segments de même longueur.}

\casecocher{Dans un triangle rectangle,
on peut utiliser le théorème de Pythagore.}

\casecocher{L'aire d'un triangle vaut
\(\dfrac{\text{base}\times\text{hauteur}}{2}\).}

\end{minipage}
\hfill
\begin{minipage}[t]{0.49\linewidth}

\casecocher{Une longueur peut être comparée à une autre
pour savoir si elle est suffisante.}

\casecocher{Un triangle isocèle possède deux côtés de même longueur.}

\casecocher{Le périmètre d'un rectangle dépend
de sa longueur et de sa largeur.}

\end{minipage}

\end{connaissancesbox}

\newpage

% =========================================================
% CORRECTION 5
% =========================================================

\pagecorrection{Problème 5 -- Le toit de l'abri}{

Comme \(M\) est le milieu de \([AB]\) :

\[
AM=\frac{AB}{2}
\]

\[
AM=\frac{8}{2}
\]

\[
AM=4\text{ m}.
\]

Le sommet \(S\) est situé à la verticale de \(M\).

Le triangle \(AMS\) est donc rectangle en \(M\).

D'après le théorème de Pythagore :

\[
AS^2=AM^2+MS^2
\]

\[
AS^2=4^2+3^2
\]

\[
AS^2=16+9
\]

\[
AS^2=25
\]

\[
AS=\sqrt{25}
\]

\[
AS=5\text{ m}.
\]

Un versant du toit mesure donc \(5\) m.

Une plaque mesure \(4{,}8\) m.

Or :

\[
4{,}8<5.
\]

\[
\boxed{\text{Une plaque de }4{,}8\text{ m n'est pas assez longue}.}
\]
}

\newpage

% =========================================================
% PROBLÈME 6 : JARDIN
% =========================================================

\titreprobleme{6}{Le jardin à engazonner}

\begin{enoncebox}

\begin{minipage}[c]{0.64\linewidth}

Un jardin est constitué d'un rectangle
et d'un triangle rectangle accolés.

Le rectangle mesure \(12\) m sur \(8\) m.

Dans le triangle rectangle, un côté mesure \(8\) m
et l'hypoténuse mesure \(10\) m.

Un sac permet d'engazonner \(35\text{ m}^2\).

\smallskip

\textbf{Combien de sacs faut-il acheter pour engazonner tout le jardin ?}

\end{minipage}
\hfill
\begin{minipage}[c]{0.30\linewidth}

\centering

\begin{tikzpicture}[scale=0.44]

\coordinate (A) at (0,0);
\coordinate (B) at (4.5,0);
\coordinate (C) at (4.5,3);
\coordinate (D) at (0,3);
\coordinate (E) at (6.75,0);

\draw[thick]
(A)--(B)--(E)--(C)--(D)--cycle;

\draw[thick]
(B)--(C);

\draw[thick]
(B) ++(0,0.32)
-- ++(0.32,0)
-- ++(0,-0.32);

\node[below left] at (A) {$A$};
\node[below] at (B) {$B$};
\node[above] at (C) {$C$};
\node[above left] at (D) {$D$};
\node[below right] at (E) {$E$};

\node[below] at ($(A)!0.5!(B)$)
{$12$ m};

\node[left] at ($(A)!0.5!(D)$)
{$8$ m};

\node[right] at ($(B)!0.5!(C)$)
{$8$ m};

\node[above right] at ($(C)!0.5!(E)$)
{$10$ m};

\end{tikzpicture}

\end{minipage}

\end{enoncebox}

\vspace{0.8mm}

\zonestravail

\vspace{0.8mm}

\begin{connaissancesbox}

{\sffamily\bfseries\color{SavoirC}
Point connaissances -- à consulter si besoin}

\smallskip

Parmi ces connaissances, lesquelles peuvent t'aider ?
Coche celles que tu utilises.

\smallskip

\begin{minipage}[t]{0.49\linewidth}

\casecocher{Dans un triangle rectangle,
on peut utiliser le théorème de Pythagore.}

\casecocher{L'aire d'un rectangle vaut
longueur \(\times\) largeur.}

\casecocher{L'aire d'un triangle vaut
\(\dfrac{\text{base}\times\text{hauteur}}{2}\).}

\end{minipage}
\hfill
\begin{minipage}[t]{0.49\linewidth}

\casecocher{Le périmètre d'une figure est la somme
des longueurs de ses côtés.}

\casecocher{Le nombre de sacs peut être obtenu
à partir de l'aire totale du jardin.}

\casecocher{Il faut acheter un nombre entier de sacs
suffisant pour couvrir toute la surface.}

\end{minipage}

\end{connaissancesbox}

\newpage

% =========================================================
% CORRECTION 6
% =========================================================

\pagecorrection{Problème 6 -- Le jardin à engazonner}{

Le triangle \(BCE\) est rectangle en \(B\).

D'après le théorème de Pythagore :

\[
CE^2=BC^2+BE^2
\]

\[
10^2=8^2+BE^2
\]

\[
100=64+BE^2
\]

\[
BE^2=100-64
\]

\[
BE^2=36
\]

\[
BE=\sqrt{36}
\]

\[
BE=6\text{ m}.
\]

L'aire du rectangle \(ABCD\) vaut :

\[
12\times8=96\text{ m}^2.
\]

L'aire du triangle rectangle \(BCE\) vaut :

\[
\frac{BE\times BC}{2}
\]

\[
\frac{6\times8}{2}
=24\text{ m}^2.
\]

L'aire totale du jardin vaut :

\[
96+24=120\text{ m}^2.
\]

Le nombre de sacs nécessaire vaut :

\[
120\div35\approx3{,}43.
\]

Trois sacs ne suffisent donc pas.

\[
\boxed{\text{Il faut acheter }4\text{ sacs}.}
\]
}

\end{document}